\section{Method}
\label{sec:method}

We present Sa2VA-Dev, which introduces a \emph{VP-aware selection bottleneck} in Sa2VA.
Our goal is to reduce the number of visual tokens consumed by the large language model while maintaining (or improving) dense grounded outputs such as referring segmentation; we instantiate the bottleneck using LIS/DiffTopK/CAS following prior differentiable selection designs.


\subsection{Overview}
\label{sec:method_overview}
Sa2VA-Dev targets a setting that is not covered by typical VQA-oriented pruning: \emph{dense grounded multi-tasking} where the same model must support (i) pixel-level mask prediction and (ii) language generation, and where additional \emph{visual prompts} (points/boxes/scribbles) inject task-critical evidence. We therefore compress the \emph{fused} image/video+VP token stream, rather than image tokens alone, so the selector can directly attend to VP-conditioned cues that drive segmentation.
Our key design is a \textbf{VP-aware selection bottleneck}: we apply budgeted token selection on the \emph{unified} visual token sequence obtained \emph{after} fusing image/video tokens with VP-derived tokens, but \emph{before} the sequence is consumed by the LLM.
This ``append-then-select'' interface makes the selector explicitly conditioned on VP information and allows the downstream LLM to operate on a shorter, VP-aligned context while preserving Sa2VA's original task interface.

Following the baseline abstraction, for an instance $i$ with input text tokens $T_i$, image/video input (denoted as $I_i$ or $V_i$), and additional visual prompt tokens $VP_i$, the model outputs answer tokens $T_o$ and optionally masks (or masklets) $M_o$:
\begin{equation}
	T_o,\ M_o = \mathrm{LLM}(\{I_i, V_i, VP_i\},\ T_i).
	\label{eq:sa2va_unified_task}
\end{equation}
Here $V_i$ denotes the raw video input, while bold symbols (e.g., $\mathbf{V}$) denote token features.
For chat-only tasks, $M_o$ is absent; for referring segmentation tasks, $M_o$ is an image mask or a spatio-temporal masklet.

We instantiate the bottleneck with a lightweight importance scorer and a train-time differentiable Top-$k$ surrogate (used only for optimization), and switch to standard hard Top-$k$ pruning at inference for real compute/memory savings.
A curriculum-weighted constraint is used to reduce the mismatch between soft gating during training and discrete pruning at inference under Sa2VA's multi-task supervision.

Figure~\ref{fig:framework} provides an overview of the end-to-end pipeline and the insertion point of the selection bottleneck.

\paragraph{Notation.}
Let $\mathbf{V}\in\mathbb{R}^{N\times D}$ denote the unified visual token features required by the LLM, where $N$ is the number of tokens and $D$ is the feature dimension.
We use a retention ratio $\rho\in(0,1]$ and set the token budget as
\begin{equation}
	k = \lfloor \rho N \rfloor.
	\label{eq:method_token_budget}
\end{equation}

\subsection{VP-to-feature conversion and unified visual feature stream}
\label{sec:method_vp}
Sa2VA supports visual prompts such as points, boxes, and scribbles.
We first rasterize VP into a binary mask $\mathbf{M}\in\{0,1\}^{H\times W}$ in the original resolution.
We then downsample the image (or each video frame) into a thumbnail view and apply an element-wise mask multiplication:
\begin{equation}
\begin{aligned}
	\bar{\mathcal{I}} &= \mathrm{Down}(\mathcal{I}) \in \mathbb{R}^{\bar{H}\times\bar{W}\times 3},\\
	\bar{\mathbf{M}} &= \mathrm{Down}(\mathbf{M}) \in \{0,1\}^{\bar{H}\times\bar{W}},\\
	\mathcal{I}^{\mathrm{vp}} &= \bar{\mathcal{I}} \odot \bar{\mathbf{M}},
\end{aligned}
\label{eq:method_vp_masked_thumb}
\end{equation}
where $\odot$ denotes broadcasting element-wise multiplication.

We encode the masked thumbnail into VP-conditioned tokens and append them to the original Sa2VA visual tokens:
\begin{equation}
\begin{aligned}
	\mathbf{X}^{\mathrm{vis}} &= \mathrm{Enc}(\mathcal{I}\ \text{or}\ \mathcal{V}) \in \mathbb{R}^{T_{\mathrm{vis}}\times d},\\
	\mathbf{X}^{\mathrm{vp}} &= \mathrm{Enc}_{\mathrm{vp}}(\mathcal{I}^{\mathrm{vp}}) \in \mathbb{R}^{T_{\mathrm{vp}}\times d},\\
	\mathbf{X} &= [\mathbf{X}^{\mathrm{vis}};\ \mathbf{X}^{\mathrm{vp}}] \in \mathbb{R}^{(T_{\mathrm{vis}}+T_{\mathrm{vp}})\times d}.
\end{aligned}
\label{eq:method_unified_tokens}
\end{equation}
For clarity, we denote this fused token stream as $\mathbf{V}$ when describing the token selection modules.
This ``append-then-select'' design is the key fusion characteristic of Sa2VA-Dev: LIS/DiffTopK operates on the unified image/video/VP token space, ensuring VP-relevant evidence is explicitly represented before pruning.

\subsection{Learnable importance scoring (LIS)}
\label{sec:method_lis}
Given the fused visual tokens $\mathbf{V}\in\mathbb{R}^{N\times D}$, LIS predicts an importance score $s_i$ for each token $i$.
Following the simplified self-attention scoring form in VisionSelector, we compute
\begin{equation}
	\mathbf{Q}=\mathbf{V}\mathbf{W}_{q},\quad \mathbf{K}=\mathbf{V}\mathbf{W}_{k},
	\label{eq:method_lis_qk}
\end{equation}
where $\mathbf{W}_{q},\mathbf{W}_{k}\in\mathbb{R}^{D\times d}$ are learnable projections and $d$ is the hidden dimension.
We then form a self-attention score matrix and obtain per-token importance by averaging interactions:
\begin{equation}
	\mathbf{A}=\frac{\mathbf{Q}\mathbf{K}^{\top}}{\sqrt{d}},\quad s_i=\frac{1}{N}\sum_{j=1}^{N}A_{ij}.
	\label{eq:method_lis_score}
\end{equation}
This scoring encourages tokens that interact strongly with many others (i.e., capturing global context) to receive higher importance.

\subsection{Train-time differentiable selection via DiffTopK}
\label{sec:method_difftopk}
The standard Top-$k$ operator is discrete and non-differentiable.
During training, we adopt a continuous relaxation to obtain a soft selection mask $\mathbf{M}_{\mathrm{soft}}\in(0,1)^{N}$:
\begin{equation}
	\mathbf{M}_{\mathrm{soft}} = \mathrm{DiffTopK}(\mathbf{s}) = \sigma(\mathbf{s}+t),\quad \sum_{i=1}^{N} M_i \approx k,
	\label{eq:method_difftopk_forward}
\end{equation}
where $\sigma(\cdot)$ is the sigmoid function and $t$ is a scalar threshold solved by bisection in the forward pass to satisfy the budget.
The soft mask can be interpreted as per-token selection probabilities and inherits the monotonicity of the sigmoid, i.e.,
\begin{equation}
	[\mathrm{DiffTopK}(\mathbf{s})]_i > [\mathrm{DiffTopK}(\mathbf{s})]_j \Leftrightarrow s_i > s_j.
	\label{eq:method_difftopk_mono}
\end{equation}
We apply the soft mask by element-wise multiplication to suppress non-critical tokens:
\begin{equation}
	\mathbf{V}_{\mathrm{pruned}} = \mathbf{M}_{\mathrm{soft}} \odot \mathbf{V}.
	\label{eq:method_pruned_tokens}
\end{equation}
This enables end-to-end gradient flow from downstream tasks to the LIS module without backpropagating through the bisection search.

\paragraph{Backward pass via implicit differentiation.}
We differentiate $\mathbf{M}=\sigma(\mathbf{s}+t)$ under the implicit constraint $\sum_{i=1}^{N}\sigma(s_i+t)=k$.
Let $\mathbf{v}=\sigma'(\mathbf{s}+t)$ with $v_i=M_i(1-M_i)$.
From the constraint, we obtain
\begin{equation}
	\frac{\partial t}{\partial s_j}=-\frac{v_j}{\sum_i v_i}.
	\label{eq:method_implicit_dt}
\end{equation}
The Jacobian of $\mathbf{M}$ with respect to $\mathbf{s}$ can be written as
\begin{equation}
	\frac{\partial \mathbf{M}}{\partial \mathbf{s}}=\mathrm{diag}(\mathbf{v})-\frac{\mathbf{v}\mathbf{v}^{\top}}{\sum_i v_i}.
	\label{eq:method_implicit_jac}
\end{equation}
Therefore, for an upstream gradient $\mathbf{g}=\frac{\partial \mathcal{L}}{\partial \mathbf{M}}$, the gradient with respect to scores is
\begin{equation}
	\frac{\partial \mathcal{L}}{\partial \mathbf{s}} = \mathbf{v}\odot\mathbf{g} - \frac{\mathbf{v}^{\top}\mathbf{g}}{\sum_i v_i}\,\mathbf{v}.
	\label{eq:method_implicit_grad}
\end{equation}

\subsection{Inference-time hard Top-$k$ pruning}
\label{sec:method_topk_infer}
At inference, we apply the standard hard Top-$k$ to obtain a binary mask $\mathbf{M}_{\mathrm{hard}}$:
\begin{equation}
    \mathbf{M}_{\mathrm{hard}}=\mathrm{TopK}(\mathbf{s}),
    \label{eq:method_hard_topk}
\end{equation}
and prune tokens accordingly (equivalently, selecting indices with ones in $\mathbf{M}_{\mathrm{hard}}$).
We evaluate multiple retention ratios (e.g., top 30\%/20\%/10\%) to quantify the trade-off between token budget and task performance, while keeping the Sa2VA baseline interface in \eqref{eq:sa2va_unified_task} unchanged.

\subsection{Curriculum annealing strategy (CAS) for multi-task training}
\label{sec:method_cas}
To bridge the discrepancy between soft selection during training and hard selection during inference, we adopt a composite objective with curriculum-based weighting.
We write the total loss as
\begin{equation}
	\mathcal{L}_{\mathrm{total}} = \mathcal{L}_{\mathrm{task}} + \lambda_t\,\mathcal{L}_{\mathrm{constraint}},
	\label{eq:method_total_loss}
\end{equation}
where $\mathcal{L}_{\mathrm{task}}$ aggregates downstream losses (segmentation, and conversation if enabled), and $\mathcal{L}_{\mathrm{constraint}}$ regularizes the selection process.

We encourage the soft mask $\mathbf{M}_{\mathrm{soft}}$ to approximate the ideal hard mask $\mathbf{M}_{\mathrm{hard}}$ via binary cross-entropy:
\begin{equation}
	\mathbf{M}_{\mathrm{hard}} = \mathrm{TopK}(\mathbf{s}),\qquad
	\mathcal{L}_{\mathrm{constraint}} = \mathrm{BCE}(\mathbf{M}_{\mathrm{soft}},\mathbf{M}_{\mathrm{hard}}).
	\label{eq:method_constraint_bce}
\end{equation}
This term pushes $\mathbf{M}_{\mathrm{soft}}$ toward polarized values (near 0 or 1), making the learned score distribution more directly applicable for hard Top-$k$ pruning.

The weighting coefficient $\lambda_t$ follows a curriculum annealing schedule for the LIS-induced selection loss. Early in training, we set $\lambda_t$ small to prioritize $\mathcal{L}_{\mathrm{task}}$; it then increases linearly to a predefined maximum:
\begin{equation}
	\lambda_t = \lambda_{\mathrm{start}} + (\lambda_{\mathrm{end}}-\lambda_{\mathrm{start}})\times \min\Big(\frac{t_{\mathrm{current}}}{t_{\mathrm{total}}},1.0\Big).
	\label{eq:method_cas_schedule}
\end{equation}

\subsection{Positional encoding under pruning: is RoPE a good fit?}
\label{sec:method_rope}
We discuss RoPE here as a \emph{design consideration} for making hard pruning behave consistently with train-time soft gating.
RoPE applies a position-dependent rotation to the query/key representations, enabling relative-position-aware attention.

Let $p_t$ denote the positional index (or a tuple of indices) associated with token $\mathbf{x}_t$.
RoPE can be written as applying a rotation operator $\mathbf{R}(p_t)$:
\begin{equation}
	\mathrm{RoPE}(\mathbf{u},p_t) = \mathbf{R}(p_t)\,\mathbf{u}.
	\label{eq:method_rope_def}
\end{equation}
For a transformer layer that computes attention over token embeddings $\mathbf{h}_t$, RoPE is typically applied to queries/keys:
\begin{equation}
\begin{aligned}
	\mathbf{q}_t &= \mathbf{R}(p_t)\,\mathbf{W}_q\mathbf{h}_t,\\
	\mathbf{k}_t &= \mathbf{R}(p_t)\,\mathbf{W}_k\mathbf{h}_t,\\
	a_{t,j} &= \mathrm{softmax}_j\Big(\frac{\mathbf{q}_t^{\top}\mathbf{k}_j}{\sqrt{d}}\Big).
\end{aligned}
\label{eq:method_rope_attention}
\end{equation}

\paragraph{What RoPE can help with.}
When inference uses hard Top-$k$, the sequence length and token subset differ from those seen in training.
If the surviving tokens keep their original positional indices (rather than being re-indexed densely after pruning), RoPE can help preserve consistent positional semantics across different token subsets.
This can reduce sensitivity to varying token counts, which is desirable when evaluating multiple retention ratios.

\paragraph{What RoPE cannot solve.}
RoPE does not address the core source of non-differentiability introduced by hard Top-$k$.
Differentiability is handled by DiffTopK during training; RoPE is largely orthogonal to this and should be viewed as a compatibility choice in the underlying MLLM rather than a replacement for DiffTopK.

\paragraph{Practical requirements in our setting.}
To make any RoPE-based claim precise, one must specify: (i) whether RoPE is 1D (sequence) or extended to 2D/spatiotemporal positions, (ii) how appended VP features are assigned positions, and (iii) whether positional indices are preserved for selected tokens at inference.
Based on these constraints, we treat RoPE as a reasonable positional encoding choice \emph{if already used by the base MLLM}, but we do not rely on RoPE to justify equivalence between near-zero soft gates and discrete token dropping.
